Agradeço inicialmente a Deus pela benção da vida, da saúde, da família e por tudo que proporcionou à minha trajetória até aqui e que eu sei que vai proporcionar no futuro.

Agradeço à minha mãe, Maria Taís Mateus Carvalho Sarmento, por toda educação, orientação, paciência, carinho, amor e compreensão que permitiu que eu chegasse até aqui. Agradeço ao meu pai, Alan Kardec Carvalho Sarmento, que também sempre foi amoroso, paciente e cujas conversas e orientações foram de caráter indispensável e incomensurável para a conclusão desse processo. Juntos me ensinaram, entre as diversas lições, o valor da educação na construção de caráter e de um espírito mais humano. Juntos são também a grande fonte de inspiração para continuar seguindo 

Agradeço também aos meus irmãos, Yule Mateus Carvalho Sarmento e Kalil Mateus Carvalho Sarmento, os quais mostraram ao longo dos anos que com companheirismo e união é possível alcançarmos juntos os mais altos sonhos.

Nesse momento quero agradecer também à minha namorada, Clara Harijan Nirvana Sousa Gomes. A sua chegada, ressaltada ainda mais pelo momento em que aconteceu, significou uma profunda transformação em minha vida. O companheirismo, o amor, a confidência, o humor e a sua força de vontade fazem com que eu a admire cada dia mais.

Agradeço também ao meu professor e orientador Otacílio da Mota Almeida sem o qual não seria possível sequer conceber esse trabalho, quanto mais concluir. Seu ensino e orientações foram tranformadores.

Agradeço aos amigos que estiveram presentes nessa jornada e aos que se juntaram no dercorrer dela. Ao amigo João Marcos Vilar o qual foi um companheiro indispensável tanto para o desenvolvimento deste trabalho, quanto para suportar as dores dessa jornado de formação que iniciamos em 2019. Aos amigos Pedro Lucas, Mateus Martini, Ryan Guedes e Guilherme Mendes que com seu humor, parceria e amor ao rock tornaram essa jornada muito mais palatável. Às amigas Laís Urano e Mirla Borges e aos amigos Bernardo Ribeiro e Geraldo Netto que mostraram, em meio às nossas maravilhosas partidas de RPG e à vivência de várias experiências, que quem tem amigos nunca está só, tem ali outra família. À amiga Izadora a qual mostrou em nossas variadas e profundas conversas as diversas maneiras de se enxergar o mundo. São muitas as pessoas importantes nesse processo e não será possível nomear todas, mas quero que todos se sintam contemplados a partir dos nomes aqui citados.

Agradeço também ao grupo PET Potência que me ampliou e completou consideravelmente minha formação enquanto profissional de engenharia e pessoa humana.

Agradeço ao grupo GRASI e aos seus componentesque forneceram estrutura e apoio para desenvolvimento deste trabalho e dos diversos trabalhos de iniciação científica produzidos ao longo do curso.

Agradeço também à Universidade Federal do Piauí por todo apoio fornecido.